%% =====================================================================
%%  PGSG manuscript -- Elsevier (elsarticle)
%%
%%  Drafting mode: single column, 12pt. For the final two-column form,
%%  swap the \documentclass line for:
%%      \documentclass[final,5p,times,twocolumn,authoryear]{elsarticle}
%%  and re-check that Table 3 still fits (it will need \small or a
%%  table* float).
%%
%%  Sections written: Introduction, Related Work, Data and Pre-processing,
%%  Method. Everything after that is a skeleton.
%%
%%  Every [TODO] marks a number, citation or figure that is deliberately
%%  left empty. None of them has been filled from memory or inference.
%% =====================================================================

\documentclass[preprint,12pt,authoryear]{elsarticle}

\usepackage{amsmath}
\usepackage{amssymb}
\usepackage{graphicx}
\usepackage{booktabs}
\usepackage{url}
\usepackage[protrusion=true,expansion=false]{microtype}
\emergencystretch=2em

%% Figure files are supplied separately; the author sets the directory, e.g.
%% \graphicspath{{figures/}}

%% \journal{[TODO: target journal]}

\begin{document}

\begin{frontmatter}

\title{Physics-Guided Scale Gating for Floe-Size Segmentation in
Sentinel-1 EW SAR Imagery}

\author[zju]{Wannian Li}
\author[zju]{Yudi Wang}
\author[zju]{Huanhao Liao}
\author[zju]{Yifang Luo}
\author[zju]{Hang Diao}
\author[zju]{Jiale Chen}
\author[zju]{Zixiong Ma}
\author[zju]{Ziyi Zhu}
\author[zju]{Yihe Wang}
\author[zju]{Fan Zhang\corref{cor1}}
\ead{f.zhang@zju.edu.cn}
\cortext[cor1]{Corresponding author.}

\affiliation[zju]{organization={State Key Laboratory of Ocean Sensing \& Ocean
                                College, Zhejiang University},
                  city={Zhoushan},
                  postcode={316021},
                  state={Zhejiang},
                  country={China}}

\begin{abstract}
[TODO: abstract. Best written last, since its strength of claim depends on the
five-fold paired result. It must stay consistent with the narrative of the
Results section: a consistent but moderate overall gain, the largest
single-class gain on big floe, no improvement on medium floe, and a readout
showing that the evidence contribution to the gate decays over training.]
\end{abstract}

\begin{keyword}
sea ice \sep floe size distribution \sep synthetic aperture radar \sep
semantic segmentation \sep physics-guided learning \sep cross-scene generalization
\end{keyword}

\end{frontmatter}

%% =====================================================================
\section{Introduction}
\label{sec:intro}
%% =====================================================================

The sea-ice Flee Size Distribution (FSD) describes how an ice cover is partitioned among floes of different horizontal extent. It enters sea-ice dynamics and thermodynamics through lateral melt, floe--floe interaction and
wave--ice coupling, and prognostic FSD schemes are now being incorporated into coupled models \citep{rothrock1984measuring,horvat2015prognostic,bateson2022floe}.
Observational FSD records remain sparse, assembled largely from high-resolution imagery that is localised in space and time \citep{toyota2006characteristics,stern2018seasonal,stern2018reconciling}, which limits the evaluation of those schemes \citep{horvat2019estimating}. Synthetic aperture radar (SAR) images the polar oceans independently of cloud cover and
illumination, and Sentinel-1 provides routine dual-polarisation coverage of the Arctic marginal ice zone \citep{torres2012sentinel1}. Producing per-pixel floe size maps from SAR is therefore a necessary component of automated operational ice charting, in which floe size is one of the coded attributes of the SIGRID-3 standard \citep{jcomm2010sigrid3,stokholm2024autoice}.

Floe size classes differ from most semantic categories in one respect that shapes everything downstream: they are defined purely by physical scale. We follow the operational coding and use six classes --- open water; cake ice($<20$~m); small floe (20--100~m); medium floe (100--500~m); big floe(0.5--2~km); and vast floe (2--10~km). At the 80~m output pixel spacing used here for Sentinel-1 Extra Wide (EW)
swath GRDM imagery, an individual small or medium floe spans
only about one to six pixels (Fig.~\ref{fig:task}). A pixel-level label is thus not a statement about brightness, nor about the semantic appearance of an object; it is a statement
about the spatial scale of the local texture. The task is, at bottom,
\emph{per-pixel scale inference}: estimating the characteristic scale of the neighbourhood around each pixel and mapping it onto an interval. General-purpose segmentation networks treat scale as one implicit property among many, learned inside features that are also carrying appearance and context. That is a structural mismatch with a task whose label set \emph{is} a partition of the scale axis.

\begin{figure}[t]
\centering
\includegraphics[width=\linewidth]{fig1_task.pdf}
\caption{The task. (a) A $41\times41$~km Sentinel-1 EW HH subset of the marginal ice zone.
(b) The corresponding floe size label; labels are constant within an ice-chart polygon, so the
supervision is regional rather than per-floe. (c) The six classes partition a physical scale
axis; the dashed line marks the processed 80~m pixel spacing, which falls inside the small-floe interval,
so individual small and medium floes span only about one to six pixels.}
\label{fig:task}
\end{figure}

To locate where a strong general-purpose model fails here, we ran a set of
diagnostic probes on a SegFormer-B2 baseline. Unless stated otherwise, every
model-side number in this paper is measured on fold~0 of the five-fold
cross-validation described in Section~\ref{sec:splits}. Two failure modes stand
out. The first concerns the \emph{quality} of the evidence. Medium floe reaches
an IoU of 0.6693 when the retained baseline checkpoint is evaluated on all
334 training scenes, compared with approximately 0.229 on the 83 held-out
validation scenes, a gap of approximately 0.440. These are scene-level
evaluation results using the same checkpoint, rather than training-batch
metrics. This places the failure in
cross-scene generalization rather than in model capacity --- consistent with
evidence that high-scoring sea-ice charting models can lean on scene- or
location-specific cues and transfer poorly to unseen regions
\citep{jalayer2025enhancing}. At the same time, the absolute radiometric level
varies widely between scenes: over our 417 scenes the standard deviation of the
per-scene mean HH backscatter is approximately 3~dB. The baseline's scale
reading is sensitive to that level. A uniform shift of about 2~dB applied to a
scene lowers hard-ice mIoU by 0.04--0.05, an order of magnitude more than the
0.005 difference between two random seeds of the same baseline. On individual
scenes the effect can be severe: Fig.~\ref{fig:radiometric} shows a fold-0
validation scene in which darkening the input by about 4~dB causes a tongue of
medium floe to be read as open water, taking the medium-floe IoU for that scene
from 0.76 to 0.03. Here and throughout, \emph{hard-ice mIoU}
denotes the mean of the small-floe and medium-floe IoU; we introduce it as a
diagnostic quantity for the two hard ice classes in this work, not as a
standard metric of the field. The picture these probes give is that the scale
evidence is contaminated by radiometric conditions.

\begin{figure}[t]
\centering
\includegraphics[width=\linewidth]{fig2_radiometric.pdf}
\caption{Scale reading is entangled with the absolute radiometric level. All four panels show
the same scene and the same grey-scale backdrop; only a constant was added to the model input.
(a) Reference labels, dominated by a tongue of medium floe. (b) With the input darkened by
0.7~z-units (about 4~dB), the tongue is read as open water and the medium-floe IoU for this
scene falls from 0.76 to 0.03. (c) Unmodified input. (d) Brightened by the same amount.
Nothing about the ice has changed.}
\label{fig:radiometric}
\end{figure}

The second failure mode concerns the \emph{structure} of the decision. The
baseline is markedly weaker on big floe than on either open water or vast floe
(IoU 0.650 against 0.969 and 0.870 respectively, fold~0). Only 65.8\% of
big-floe pixels are recovered: 16.7\% are read as medium floe and 1.5\% as small
floe, while a further 15.5\% are absorbed into vast floe. The errors therefore
run in both directions along the scale axis --- part of a large floe is read as a
field of smaller ones, and part of it is merged into the class above. A pixel
deep inside a big floe has, at fine scale, much the local texture of a field of
small floes, and at coarse scale little to distinguish it from a vast floe;
placing it correctly requires the two scales to be read against one another.
When each scale branch is weighted independently and nothing couples them,
neither comparison is available. Taken together, the two failure modes
point at a single structural deficiency: the evidence on which the per-pixel
scale decision rests is both \emph{contaminated}, by scene radiometry, and
\emph{isolated}, by the absence of cross-scale coordination.

We address this with a single design principle, which we call
\emph{evidence--decision decoupling}: quantities that are physically grounded
and stable under changes in scene radiometry serve as scale evidence; the
per-pixel scale decision is conditioned on that evidence; and the evidence
enters the network only as a conditioning variable, never in parallel with the
main feature path. We instantiate it with three mechanisms. (i)
\emph{Contrast-invariant texture evidence}: a differentiable variogram profile
which, after normalisation across lags, is invariant to affine radiometric
transformations $f \mapsto af+b$ and retains only correlation length, that is,
texture coarseness. (ii) A \emph{polarimetric prior}: the HH$-$HV difference in
dB --- a ratio in linear units --- used as a calibration-light fingerprint of
the scattering mechanism. (iii) \emph{Cross-scale communication and an
evidence-conditioned competitive scale gate}: at each pixel the four scale
features first exchange information through a single scale-axis transformer
layer, after which a budget-conserving softmax gate distributes a fixed scale
budget among them, conditioned on the evidence. The two projections through
which these modules act on the network --- the output of the scale-axis
transformer and the final layer of the gate --- are zero-initialized, so at the
start of training the forward pass is bit-for-bit identical to the baseline; whatever they contribute comes from learning and not
from a perturbation of the initialisation.

Our contributions are the following.
\begin{itemize}
\item \textbf{Problem formulation.} We formulate SAR floe size segmentation as
      per-pixel scale inference, and use diagnostic probes to identify
      radiometric contamination of the scale evidence and the absence of
      cross-scale coordination as the two principal sources of failure.
\item \textbf{Method.} We propose PGSG (Physics-Guided Scale Gating), an
      evidence--decision decoupled architecture combining contrast-invariant
      texture evidence, a polarimetric prior, and a budget-conserving,
      evidence-conditioned scale gate.
\item \textbf{Component-level attribution that can be read.} Every added module
      acts on the network through a zero-initialized projection and is
      instrumented with training-time readouts --- the
      gate budget distribution, the activation share contributed by the evidence
      at the gate input, and the residual norm ratio of the scale-axis
      transformer. Whether a module \emph{is used} and whether it \emph{helps}
      can therefore be judged separately, which is what allows us to report both
      positive and negative outcomes.
\end{itemize}
On fold~0, macro mIoU rises from 0.5031 to 0.5109. The largest single-class gain
is on big floe ($+0.026$), where the fragmentation described above is partly
repaired; the hard ice classes improve only slightly, and medium floe does not
improve at all. [TODO: five-fold paired result.]

%% =====================================================================
\section{Related Work}
\label{sec:related}
%% =====================================================================

\subsection{SAR sea-ice retrieval and ice-chart automation}

Automated interpretation of SAR for operational ice charting has been driven in
recent years by the AI4Arctic Sea Ice Challenge dataset and the AutoICE
challenge, which posed sea ice concentration (SIC), stage of development (SOD)
and floe size (FLOE) as three parallel segmentation tasks over Sentinel-1 scenes
paired with ice charts from the Greenland and Canadian ice services
\citep{stokholm2024autoice,buushinkler2023rttai4arctic}.
Submissions converged on multi-task U-Net variants with task-specific heads and
auxiliary inputs \citep{chen2024mmseaice}. Earlier deep-learning work on
Sentinel-1 sea ice concentrated on ice type classification from patch-wise
convolutional networks trained against expert charts
\citep{boulze2020classification}. A recent re-analysis of the AutoICE benchmark
used spatial cross-validation and attribution methods to show that
high-scoring solutions can rely on geospatial encoding to a degree that does not
survive transfer to regions outside the benchmark split
\citep{jalayer2025enhancing}.

FSD retrieval proper has developed along a separate track, on imagery an order
of magnitude finer than EW-mode SAR. The classical pipeline is ice--water
segmentation followed by separation of touching floes, using thresholding with
manual correction \citep{toyota2006characteristics,toyota2011size}, distance
transforms and watershed with rule-based boundary revalidation
\citep{hwang2017practical}, or active contours driven by gradient vector flow
\citep{zhang2015imageprocessing}; more recently, classical segmentations have
been used to label training data for deep instance segmentation
\citep{zhang2023icefloe}. What these approaches share is a dependence on
resolution fine enough to delineate individual floes. On the processed 80~m
grid used here for Sentinel-1 EW imagery that assumption does not hold, and per-pixel FLOE
segmentation at this resolution has attracted little dedicated methodology; the
prevailing practice is to reuse general-purpose segmentation networks.

\subsection{Semantic segmentation and multi-scale fusion}

Modern semantic segmentation descends from fully convolutional networks
\citep{long2015fcn} and encoder--decoder designs with skip connections
\citep{ronneberger2015unet}, with context aggregated by dilated convolution and
spatial pyramid pooling \citep{chen2018deeplab} or by maintaining high-resolution
representations throughout the network \citep{wang2021hrnet}. Transformer
backbones have since become standard, either through hierarchical windowed
attention \citep{liu2021swin} or through a hierarchical encoder paired with a
lightweight all-MLP decoder \citep{xie2021segformer}.

How multi-scale features are combined falls broadly into two paradigms. The
first concatenates features from several strides with fixed weights and fuses
them with a learned projection; the all-MLP head of \citet{xie2021segformer} is
a representative case. The second selects among scales adaptively, weighting
each scale per pixel through an attention mechanism
\citep{chen2016attentiontoscale}, choosing among kernels of different receptive
field \citep{li2019sknet}, or generating input-dependent convolution weights
\citep{chen2020dynamicconv}. What is common to the adaptive family is that the
input to the selection mechanism is the feature tensor itself. Whatever
confounds the features therefore also confounds the selection --- which is
precisely the gap the evidence-conditioning in this work is aimed at.

\subsection{Physical priors and invariant representations in SAR}

Polarimetric ratios and decompositions are long-established descriptors of
sea-ice scattering. Co- and cross-polarised ratios separate open water from ice
types with thresholds that are comparatively stable across acquisitions
\citep{geldsetzer2009seaice}, and the broader family of polarimetric parameters
and their sea-ice applications has been reviewed by \citet{shokr2023polarimetric}.
The incidence-angle dependence of backscatter is itself ice-type dependent and
differs between co- and cross-polarised channels, which is why explicit
normalisation schemes have been developed \citep{aldenhoff2020incidence}.

Texture statistics under speckle form a second strand. Grey-level co-occurrence
features \citep{haralick1973textural} remain the most widely used descriptors for
SAR sea ice \citep{clausi2001comparison}; notably, \citet{shokr1991evaluation}
already modified individual co-occurrence parameters so as to be invariant under
linear grey-tone transformations, which is the same concern that motivates the
texture evidence in this work. Geostatistical descriptors offer a more direct
route to spatial scale: variogram (structure function) analysis characterises
surface roughness and correlation length \citep{garrigues2007variograms}, and has
been applied to sea ice both for classifying ice provinces
\citep{herzfeld2016geostatistical} and for estimating ice concentration and floe
scale from SAR intensity through fitted stochastic models
\citep{wang2016variograms}. Scattering transforms provide a related construction
with explicit invariance guarantees \citep{bruna2013scattering}. Speckle sets a
lower bound on the scales at which any of this is meaningful, since it behaves as
multiplicative noise concentrated at the finest scales
\citep{lee1981speckle,lee1981refined,argenti2013tutorial}.

A third strand addresses generalization across acquisition conditions directly,
through adversarial alignment of domain representations \citep{ganin2016dann},
contrastive pre-training \citep{chen2020simclr}, or evaluation protocols that
respect spatial structure and thereby avoid optimistic estimates from leakage
between correlated samples \citep{roberts2017crossvalidation,jalayer2025enhancing}.

Across these strands, physical quantities are almost always consumed as
\emph{additional input channels}. Using one as a conditioning variable that
steers a decision mechanism, while deliberately keeping it out of the main
representation path, is much less common, and it is the position this work takes.
%% [TODO: PA^2-Swin. Add the one-sentence lineage statement here once the
%%  formal reference is available.]

%% =====================================================================
\section{Data and Pre-processing}
\label{sec:data}
%% =====================================================================

\subsection{Dataset}

The dataset comprises 417 Sentinel-1 EW GRDM scenes in dual-polarisation HH/HV
(1SDH), acquired between April 2022 and June 2026. During SNAP terrain
correction, the imagery is resampled to an output grid with 80~m pixel
spacing. This is the processed grid spacing used in this study, not the
native product pixel spacing or a claim of 80~m spatial resolution.
Training and validation use this same grid without further scene
downsampling. Sensors are Sentinel-1A (342 scenes) and Sentinel-1C (75 scenes).
Geographically the dataset is anchored in East Greenland --- Cape Farewell,
Central East and South East, with Danish Meteorological Institute regional ice
charts as the label source --- and supplemented with 19 Canadian Arctic scenes
(Beaufort Sea, Canadian Arctic Archipelago, Hudson Bay) carrying ice charts of
the AI4Arctic lineage, added to broaden the coverage of the large floe classes.
%% [TODO] Insert the per-region breakdown of labelled pixels by class here. Do
%% NOT assert that the big and vast floe classes come predominantly from the
%% Canadian Arctic: with 19 of 417 scenes from that region, the claim is not
%% supported by anything measured so far and must wait for the breakdown.
Seasonally the dataset is uneven: 59 winter, 187 spring, 159 summer and
12 autumn scenes. Table~\ref{tab:dataset} summarises the composition. The
dataset was assembled for this work and is not a public release.

\begin{table}[t]
\centering
\caption{Dataset composition. Class pixel shares are computed over labelled
pixels only, excluding the ignore class. The regional split of the class shares
is not yet quantified and is therefore not asserted in the text.}
\label{tab:dataset}
\begin{tabular}{llr}
\toprule
Attribute & Value & Count \\
\midrule
Scenes        & total                            & 417 \\
Sensor        & Sentinel-1A                      & 342 \\
              & Sentinel-1C                      & 75 \\
Region        & East Greenland                   & 398 \\
              & Canadian Arctic                  & 19 \\
Season        & winter                           & 59 \\
              & spring                           & 187 \\
              & summer                           & 159 \\
              & autumn                           & 12 \\
\midrule
Class share   & open water                       & [TODO] \\
              & cake ice                         & [TODO] \\
              & small floe                       & [TODO] \\
              & medium floe                      & [TODO] \\
              & big floe                         & [TODO] \\
              & vast floe                        & [TODO] \\
\bottomrule
\end{tabular}
\end{table}

\subsection{Labels}

Labels are derived from SIGRID-3 ice-chart polygons \citep{jcomm2010sigrid3}
rasterised onto the SAR grid, with the official floe-size codes mapped onto the
six classes of Section~\ref{sec:intro}. Each polygon is assigned its dominant
floe form, accepted when total concentration is at least 60\% and the dominant
form accounts for at least 65\% of the ice; these thresholds are looser than is
customary, chosen to retain recall on small and medium floes inside the marginal
ice zone. Land, SAR border regions, ambiguous polygons and fast ice are assigned
the ignore index 255 and are excluded from both loss and evaluation.

One property of this supervision governs much of the method that follows: the
label is constant across a polygon. It is regional supervision, not
floe-by-floe delineation, and it carries no information about individual floe
boundaries. What a model can learn from it is the texture scale characteristic
of a region, not the detection of single floes --- which is also why the
evidence constructed in Section~\ref{sec:method} takes the form of regional
statistics rather than instance-level features.

\subsection{SNAP processing chain}

Scenes were processed with the Microwave Toolbox (version 12.0.1) of ESA SNAP
\citep{esa2025snap} in the following order.
\begin{enumerate}
\item \textbf{Orbit correction.} Sentinel Precise orbit state vectors,
      automatically downloaded, polynomial degree 3.
\item \textbf{Thermal noise removal.}
\item \textbf{GRD border noise removal.} Border limit 500, trim threshold 0.5.
\item \textbf{Radiometric calibration} to $\sigma^0$ in linear units.
\item \textbf{Speckle filtering.} Refined Lee \citep{lee1981refined}, filter
      size $3\times3$, window $7\times7$, target window $3\times3$, with the
      equivalent number of looks estimated from the data.
\item \textbf{Terrain correction.} Copernicus 90~m Global DEM (COP-DEM GLO-90)
      \citep{opentopography2021copdem}, bilinear resampling of both DEM and
      image, output pixel spacing 80~m. The ellipsoidal incidence angle is
      retained as an output band.
\item \textbf{Conversion to dB.}
\item \textbf{Invalid pixels} set to $-9999$.
\item \textbf{GLCM texture} \citep{haralick1973textural} computed on HH:
      probabilistic quantiser, 64 grey levels, $11\times11$ window,
      displacement 2, averaged over all four directions; contrast,
      dissimilarity and homogeneity are retained.
\end{enumerate}
The incidence angle band used downstream is the \emph{ellipsoidal} incidence
angle, not the local incidence angle.

\subsection{Input channels and normalisation}
\label{sec:channels}

The network input has seven channels: HH and HV in dB, clipped to
$[-35, 5]$; the ellipsoidal incidence angle; the three GLCM texture bands; and
the polarisation ratio HH$-$HV, derived online. Continuous channels are
standardised with a single global mean and standard deviation per channel,
computed once over the whole archive; invalid pixels are set to zero after
standardisation.

We use global rather than per-scene standardisation by design, not as a result
of an experimental comparison. The absolute radiometric level carries genuine
discriminative information --- open water and ice differ in mean backscatter ---
and per-scene standardisation would remove that information along with the
nuisance variation. The strategy pursued here is instead to keep the absolute
level available to the main path while supplying the decision mechanism with
evidence that is invariant to it.

\subsection{Sampling, splits and evaluation protocol}
\label{sec:splits}

\paragraph{Tile-weighted sampling.} Each scene is indexed into $512\times512$
tiles aligned with the storage chunking, and $256\times256$ patches are cropped
from within them. Sampling probabilities over tiles are obtained by iterative
proportional fitting \citep{deming1940leastsquares} so that the
\emph{pixel-level} class composition of a batch approaches a target mixture
(open water 0.15, cake ice 0.10, and 0.1875 for each of the four remaining ice
classes), which counteracts the strong class imbalance of the archive. Tile
histograms are weighted by the coverage kernel induced by patch cropping, so
that the fitted probabilities refer to the pixels a batch actually sees rather
than to whole-tile contents.

\paragraph{Augmentation.} [TODO: augmentation list, per the training
configuration.]

\paragraph{Cross-validation.} Scenes are split with a grouped five-fold scheme.
Groups are formed by taking the transitive closure of two relations --- same
sensor and same relative orbit with adjacent frames, and acquisition start times
within 2~hours --- which collapses the 417 scenes into 292 groups and prevents
a scene and its neighbouring frames from straddling the train/validation
boundary \citep{roberts2017crossvalidation}. Folds are then balanced on the
six-class pixel shares, the scene count and the number of Sentinel-1C scenes,
giving 83--84 scenes per fold, 15 Sentinel-1C scenes per fold, and per-class
pixel shares within 17--22\% of the archive values. The fold assignment file is
frozen and shared by all experiments reported here. The experiments reported
in the present manuscript use fold~0: 334 scenes for training and 83 disjoint
scenes for validation, together covering all 417 scenes. The validation scenes
are used for checkpoint selection and early stopping as well as for the
reported validation metrics; they are not an independent test set. Table~\ref{tab:folds}
reports the resulting balance.

\begin{table}[t]
\centering
\caption{Balance check on the five-fold grouped split.}
\label{tab:folds}
\begin{tabular}{lcccc}
\toprule
Fold & Scenes & Sentinel-1C & Groups & Six-class share range \\
\midrule
0 & 83 & 15 & [TODO] & [TODO] \\
1 & [TODO] & 15 & [TODO] & [TODO] \\
2 & [TODO] & 15 & [TODO] & [TODO] \\
3 & [TODO] & 15 & [TODO] & [TODO] \\
4 & [TODO] & 15 & [TODO] & [TODO] \\
\bottomrule
\end{tabular}
\end{table}

\paragraph{Validation and inference.} Validation uses non-overlapping
$256\times256$ tiles; only tiles containing valid labels are passed forward.
Confusion matrices are accumulated per scene and all metrics are derived from
the aggregated confusion matrix, so that metrics are independent of scene size
and are not subject to any cap on the number of pixels evaluated. The reported
metrics are macro mIoU over the six classes with index 255 ignored, macro-F1,
and the hard-ice diagnostic defined in Section~\ref{sec:intro}. The checkpoint
retained is the one maximising a moving average of mIoU over the three most
recent validations, with early stopping after 20 validations without
improvement. The moving average is part of the evaluation protocol rather than a
result: selecting on a single validation peak would bias the reported value
upward by picking a favourable fluctuation.

A limitation of the design should be stated here and is returned to in
Section~\ref{sec:limitations}. The archive is geographically concentrated: 398
of 417 scenes come from East Greenland, so any correlation between class and
region that does exist cannot be broken by the split. Grouped cross-validation
prevents leakage between a scene and its neighbouring frames, but it cannot
remove a correlation that is built into the sampling of the archive itself.
%% [TODO] Once the per-region class pixel breakdown is available, state the
%% actual degree of class/region correlation here. The earlier working
%% assumption -- small/medium from East Greenland, big/vast from the Canadian
%% Arctic -- is not currently supported and must not be reinstated without it.

%% =====================================================================
\section{Method}
\label{sec:method}
%% =====================================================================

\subsection{Problem formulation and overview}

Let $x \in \mathbb{R}^{7 \times H \times W}$ denote the input tensor of
Section~\ref{sec:channels}. The task is to predict a label
$y(u) \in \{0,\dots,5\}$ at every pixel $u$, where class $c$ corresponds to a
floe size interval $[l_c, u_c]$ on the physical scale axis. Framed as scale
inference, the model must estimate a characteristic scale of the neighbourhood
of $u$ and map it into one of those intervals.

The backbone is a MiT-B2 encoder \citep{xie2021segformer} initialised from
ImageNet \citep{deng2009imagenet}, producing features at strides
$\{4, 8, 16, 32\}$. Each is projected to 256 channels and upsampled to stride~4,
giving $S=4$ aligned scale features. These pass through a scale-axis transformer
layer, then through an evidence-conditioned scale gate, and are concatenated
(1024 channels) and fused by a $1\times1$ convolution with batch normalisation
and ReLU before the classification head.

Two evidence streams run alongside the main path. The stride-4 encoder feature
feeds the variogram module (Section~\ref{sec:variogram}), producing 64 channels;
the polarisation ratio channel of the input feeds the polarimetric branch
(Section~\ref{sec:polar}), producing 32 channels. Their concatenation is the
evidence tensor $E \in \mathbb{R}^{96 \times H/4 \times W/4}$.

The central design principle is that $E$ reaches only the decision modules ---
the evidence token of the scale-axis transformer and the input of the gate ---
and is never added to, or concatenated with, the main feature path. This is what
we mean by evidence--decision decoupling, and the constraint at the end of
Section~\ref{sec:variogram} follows from it. Figure~\ref{fig:arch} shows the two
paths in separate colours, together with the two zero-initialized projections.

\begin{figure}[t]
\centering
\includegraphics[width=\linewidth]{fig3_architecture.pdf}
\caption{PGSG. The main feature path (blue) runs from the input through the MiT-B2 encoder to
four scale features, which exchange information in a scale-axis transformer and are then
weighted by a budget-conserving softmax gate before fusion. The evidence path (red, dashed)
supplies the variogram profile and the polarimetric prior to the two decision modules only;
it never joins the main feature path. Circled zeros mark the two zero-initialized projections,
which make the forward pass at initialisation identical to the baseline.}
\label{fig:arch}
\end{figure}

\subsection{Contrast-invariant texture evidence}
\label{sec:variogram}

The stride-4 encoder feature is first projected by a $1\times1$ convolution to
$G=16$ scalar fields $f$. For each lag $r \in \{1,2,4,8,16\}$ --- corresponding
to 320~m, 640~m, 1.28~km, 2.56~km and 5.12~km on the ground --- a local
structure function is computed over an $11\times11$ window $\mathcal{W}(u)$:
\begin{equation}
\gamma_r(u) \;=\;
\frac{1}{|\mathcal{W}(u)|}\sum_{v \in \mathcal{W}(u)}
\frac{1}{4}\sum_{d \in \{\pm r e_h,\; \pm r e_w\}}
\bigl(f(v) - f(v+d)\bigr)^2 .
\label{eq:variogram}
\end{equation}
The profile is then normalised across lags and clipped,
\begin{equation}
v_r(u) \;=\; \operatorname{clamp}\!\left(
\log\bigl(\gamma_r(u) + \epsilon\bigr)
- \log\Bigl(\tfrac{1}{|R|}\textstyle\sum_{r' \in R} \gamma_{r'}(u) + \epsilon\Bigr),
\; -6,\; 6 \right),
\label{eq:vprofile}
\end{equation}
where $R$ is the set of lags and $\epsilon$ a numerical stabiliser.

\paragraph{Invariance.} Let $f \mapsto af + b$ with $a \neq 0$. Differencing
removes the offset, so $\gamma_r \mapsto a^2\gamma_r$ for every lag
simultaneously. Up to $\epsilon$, each log term in \eqref{eq:vprofile} is
shifted by $2\log|a|$, and the shift cancels against the normaliser, so
$v$ is unchanged. The profile therefore discards amplitude entirely and keeps
only shape --- where $\gamma$ approaches its sill as a function of lag, which is
a dense estimate of the correlation length. Correlation length is exactly what
distinguishes small from medium from big floe, and amplitude is exactly the
quantity that varies by several dB between scenes for reasons unrelated to floe
size. The construction is thus targeted at motivation~B of
Section~\ref{sec:intro}. The concern is not new --- \citet{shokr1991evaluation}
modified co-occurrence parameters for invariance under linear grey-tone
transformations three decades ago --- but three things here are: the estimator
is differentiable and trained end to end; it retains a \emph{profile} over lags
rather than collapsing to a scalar; and it is used to condition a decision
rather than as an input feature.

Fig.~\ref{fig:variogram} illustrates both halves of this statement: the raw
structure function of a given class is displaced by the radiometric level of the
scene, while the normalised profile is unchanged by the same transformation.

\paragraph{Design choices.} Three details are argued from first principles
rather than from an experimental comparison. Displacements are taken
symmetrically in four directions so that the estimator is equivariant under
flips and its measurement centre does not drift. A single fixed $11\times11$
window is shared by all lags, so that every lag is estimated under the same
noise budget and differences along the profile reflect the signal rather than
the estimator. The smallest lag is 320~m rather than one pixel, because speckle
behaves as multiplicative noise dominating the finest scales
\citep{lee1981speckle,argenti2013tutorial}; a one-pixel lag would measure mostly
speckle, and the texture signal would be buried.

The profile is projected to 64 channels by a $1\times1$ / $3\times3$ mixing
block. It is \emph{not} added as a residual to the main features; it serves only
as evidence for the gate and the scale-axis transformer. This follows from the
decoupling principle: placed in parallel with the main path, the evidence would
compete with the gate for the same representational bandwidth, and its role as a
clean conditioning signal would be lost.

\begin{figure}[t]
\centering
\includegraphics[width=\linewidth]{fig4_variogram.pdf}
\caption{The construction of Section~\ref{sec:variogram} applied to the $4\times$-pooled HH
channel, which shares the 320~m grid of the stride-4 encoder features on which the module
actually operates. Lines are class medians over ten near-pure tiles (nine for vast floe), each
taken from a different scene; shading is the interquartile range. (a) The raw structure
function. The spread within a class is large and the classes overlap, because the amplitude is
set mainly by the radiometric level of the scene; an affine change $f \mapsto af+b$ (coloured
dashed) displaces a profile bodily. (b) After normalisation across lags, the transformed
profiles (black dashed) coincide with the originals to $4\times10^{-5}$, so only the shape
survives. The median profiles then order monotonically with floe size---steepest for small floe
and flattest for big and vast floe, which are indistinguishable at these lags---although the
per-tile ranges still overlap.}
\label{fig:variogram}
\end{figure}

\subsection{Polarimetric prior}
\label{sec:polar}

The second evidence stream is the polarisation ratio $p = \mathrm{HH} -
\mathrm{HV}$ in dB, which is a ratio in linear units. It is average-pooled
$4\times4$ to stride~4, which suppresses residual speckle, and passed through
two $3\times3$ convolutions with batch normalisation and ReLU to give 32
channels.

Physically, $p$ responds to the difference in scattering mechanism between
specular open water, level ice and deformed ice, for which co- and
cross-polarised returns behave differently \citep{geldsetzer2009seaice,
shokr2023polarimetric}. Taking a difference in dB cancels multiplicative
common-mode terms, so the quantity is comparatively insensitive to absolute
calibration. It also attenuates, though it does not remove, the incidence-angle
dependence: the angular slope of backscatter is itself polarisation- and
ice-type-dependent, so a residual angular term survives the difference
\citep{aldenhoff2020incidence}. We therefore treat $p$ as a
calibration-light, not calibration-free, descriptor.

\subsection{Cross-scale communication}
\label{sec:xscale}

The motivation is the fragmentation described as motivation~C in
Section~\ref{sec:intro}. Under purely per-scale weighting, the fine-scale branch
has no way of knowing that the pixel it is describing lies inside a single large
floe, so the ridges and internal texture of that floe are read as a field of
small or medium floes. The remedy is to let the scales talk to each other before
they are weighted.

At every pixel, the $S=4$ scale features (256 channels each) are projected by a
linear map to $d=64$ and treated as four tokens; a fifth token is formed from
the evidence tensor $E$ by a $1\times1$ convolution to 64 channels. Learnable
scale embeddings are added, and the five tokens pass through a single pre-LN
transformer layer with four attention heads and an MLP ratio of 2. The four
scale tokens are read back and mapped by a zero-initialized linear layer
$64 \to 256$, which is added as a residual to the corresponding scale feature.
The sequence length is five, so the layer is cheap despite being applied at
every pixel. [TODO: measured parameter and FLOP increment of this layer; see
also Section~\ref{sec:training} for the whole-model figure.]

Attention is implemented with batched matrix multiplications rather than a fused
attention kernel, because one sequence per pixel produces a number of sequences
beyond the grid limits such kernels assume.

\subsection{Evidence-conditioned competitive scale gate}
\label{sec:gate}

The concatenated scale features $X$ (1024 channels) and the evidence $E$ (96
channels) are concatenated to 1120 channels and passed through a $1\times1$
convolution to 64 channels with batch normalisation and ReLU, a $3\times3$
convolution with the same width and normalisation, and a final $1\times1$
convolution producing $S=4$ logits $a_s(u)$. The gate values are
\begin{equation}
g_s(u) \;=\; S \cdot \operatorname{softmax}_s\bigl(a(u)\bigr),
\qquad \sum_{s=1}^{S} g_s(u) \equiv S ,
\label{eq:gate}
\end{equation}
so the total gate budget is conserved at every pixel. $X$ is then multiplied
block-wise by the corresponding $g_s$ before the fusion convolution.

The budget constraint is a design decision, argued as follows rather than
established experimentally. The softmax form means that raising one scale
necessarily lowers another: the module's only degree of freedom is to
\emph{choose among} scales. It cannot uniformly amplify or attenuate all scales
at once, and that excluded degree of freedom has no counterpart in the task
semantics --- a per-pixel scale inference problem has nothing to say about
overall gain --- while it would add capacity that noise could exploit.
Constraining the module to the operation the task actually calls for is the
conservative choice.

The final convolution of the gate is zero-initialized, so $a \equiv 0$ and
$g \equiv 1$ at initialisation and the gated fusion is identical to ungated
concatenation.

\subsection{Training and the zero-initialisation discipline}
\label{sec:training}

The loss is six-class cross-entropy with index 255 ignored. Optimisation uses
AdamW \citep{loshchilov2019adamw} with discriminative learning rates ---
$10^{-5}$ for the encoder, $10^{-4}$ for the decoder head and all added modules
--- mixed precision with gradient scaling, and gradient clipping at 5.0. An
epoch is 300 batches of size 16; training runs for at most 200 epochs with the
early stopping and checkpoint selection of Section~\ref{sec:splits}.
[TODO: hardware, wall-clock time per fold, total parameter count, and the
parameter and FLOP increment of the full model over the baseline, i.e.\ the
scale-axis transformer, the two evidence streams and the gate combined.]

The zero-initialisation discipline deserves a statement of its own, because the
component-level attribution claimed in Section~\ref{sec:intro} rests on it. The
final projection of the scale-axis transformer and the final convolution of the
gate --- the two points at which the added modules act on the network --- are
zero-initialized. At $t=0$ the forward pass of every variant is therefore
bit-for-bit identical to the baseline, a property verified by automated
assertions rather than by inspection. In addition, all added modules are
constructed after the encoder and the decoder head, so that under a fixed random
seed the weights shared between variants are bit-for-bit identical as well.
Together these two properties mean that the difference between any two
ablation variants originates in training dynamics alone: the single-variable
comparison holds at the level of the implementation, not merely by intent.

\subsection{Readouts for component-level attribution}
\label{sec:readouts}

To make both positive and negative outcomes decidable, four quantities are
logged during training: the distribution of the gate budget across scales and
its selection entropy; the share of the activation at the gate input
attributable to the evidence channels; the ratio of the residual norm added by
the scale-axis transformer to the norm of the scale features it is added to; and
the fraction of variogram profile entries hitting the clamp bounds of
\eqref{eq:vprofile}.

These readouts separate two questions that are otherwise easy to conflate:
whether a module \emph{is being used} by the trained network, and whether it
\emph{improves} the metric. A module can be heavily used and still not help; a
module that is not used at all cannot be blamed for a null result, and the
correct conclusion in that case concerns optimisation rather than design. The
attribution in the ablation refers back to these readouts. Their numerical
values belong to the ablation and are not reported here.

%% =====================================================================
%%  SKELETON SECTIONS -- not part of the current drafting scope
%% =====================================================================

\section{Experimental setup}
\label{sec:experiments}

Five configurations are compared. (a) The \emph{baseline} is the MiT-B2 encoder
with the projection-and-fusion decoder of Section~\ref{sec:method} and no added
modules. (b) adds the two evidence streams and the evidence-conditioned scale
gate, but not cross-scale communication. (c) is the full PGSG model, that is
(b) with the scale-axis transformer of Section~\ref{sec:xscale}. Two further
configurations are probes rather than candidate components: (d) adds to (b) an
adversarial scene-alignment objective of the domain-adversarial type
\citep{ganin2016dann}, and (e) adds to (c) a polarimetric bias attention module.
[TODO: one or two sentences of implementation detail for (d) and (e), sufficient
for a reader to reproduce the probes.]

All five share the training recipe of Section~\ref{sec:training}, the same
random seed, and the frozen fold assignment of Section~\ref{sec:splits}. Because
of the zero-initialisation discipline, the forward pass of every configuration
at $t=0$ is bit-for-bit that of the baseline, and the weights shared between
configurations are bit-for-bit identical; differences between rows therefore
originate in training dynamics rather than in initialisation.

All model-comparison results reported below are measured on the fold~0
validation set (83 scenes), using models trained on the other 334 scenes,
with metrics derived from the confusion matrix aggregated over those
validation scenes as described in Section~\ref{sec:splits}. The separate
training-scene evaluation in Section~\ref{sec:intro} uses all 334 training
scenes only to diagnose the training--validation gap; those scenes are not
included in the reported validation scores. We note the restriction explicitly
because it governs how the numbers may be read: the five-fold aggregate is a
different evaluation set, and its absolute values are not comparable with the
fold-0 values quoted below. Every comparison in this paper is therefore made
within fold~0.

\section{Results}
\label{sec:results}

\subsection{The main chain}

Table~\ref{tab:ablation} reports the ablation. Reading the main chain
incrementally is more informative than reading the endpoints alone, because the
two additions do different things and one partly compensates for the other.

Adding the evidence streams and the gate, (a)~$\rightarrow$~(b), moves macro
mIoU from 0.5031 to 0.5009 --- a change of $-0.002$, which on its own would
suggest the module does nothing useful. The aggregate hides a redistribution.
Small floe rises from 0.220 to 0.238 ($+0.018$) and the hard-ice diagnostic from
0.2243 to 0.2343 ($+0.010$), the highest hard-ice value anywhere in the table;
meanwhile big floe falls by 0.010 and vast floe by 0.021. Conditioning the
per-pixel scale decision on radiometry-invariant evidence therefore buys
accuracy on the hard classes, and pays for it on the large ones. That is
consistent with what the mechanism is built to do, and it also shows what it
does not do: conditioning alone leaves the large-floe fragmentation of
motivation~C untouched, and indeed makes it slightly worse.

Adding cross-scale communication, (b)~$\rightarrow$~(c), moves macro mIoU from
0.5009 to 0.5109 ($+0.010$). The gain is concentrated precisely on the classes
(b) had given up: big floe $+0.036$ and vast floe $+0.025$. The hard classes
return a little of their gain (small floe $-0.007$, hard-ice $-0.004$). The
pattern is what the mechanism predicts --- a pixel inside a large floe can only
be recognised as such from the coarse scale, and the scale-axis layer is what
lets that information reach the fine branch --- but we note that the pattern is
consistent with the mechanism rather than proof of it.

End to end, (a)~$\rightarrow$~(c) gives $+0.008$ macro mIoU, from 0.5031 to
0.5109. The decomposition is uneven: big floe $+0.026$, the largest single-class
gain in the study; small floe $+0.011$; cake ice $+0.007$; vast floe $+0.004$;
open water $-0.001$, which is no material degradation of an already saturated
class; and medium floe unchanged. The hard-ice diagnostic improves by $0.006$.

\begin{table}[t]
\centering
\footnotesize
\setlength{\tabcolsep}{4pt}
\caption{Ablation on fold~0 (83 validation scenes, metrics derived from the
aggregated confusion matrix). Per-class values are IoU. Hard-ice is the mean of
the small-floe and medium-floe IoU, a diagnostic quantity defined in this work.}
\label{tab:ablation}
\begin{tabular}{@{}p{0.30\linewidth}cccccccc@{}}
\toprule
 & macro & \multicolumn{6}{c}{IoU by class} & hard- \\
\cmidrule(lr){3-8}
Configuration & mIoU & water & cake & small & medium & big & vast & ice \\
\midrule
(a) Baseline
  & 0.5031 & 0.969 & 0.081 & 0.220 & 0.229 & 0.650 & 0.870 & 0.2243 \\[2pt]
(b) + evidence-conditioned scale gate
  & 0.5009 & 0.970 & 0.078 & 0.238 & 0.231 & 0.640 & 0.849 & 0.2343 \\[2pt]
\textbf{(c) PGSG (full): (b) + cross-scale communication}
  & \textbf{0.5109} & 0.968 & 0.088 & 0.231 & 0.229 & 0.676 & 0.874 & 0.2302 \\
\midrule
\multicolumn{9}{@{}l}{\emph{Additional probes}} \\
(d) (b) + adversarial scene alignment
  & 0.4815 & 0.951 & 0.067 & 0.188 & 0.251 & 0.587 & 0.845 & 0.2198 \\[2pt]
(e) (c) + polarimetric bias attention
  & 0.5103 & 0.969 & 0.088 & 0.233 & 0.206 & 0.700 & 0.866 & 0.2193 \\
\bottomrule
\end{tabular}
\end{table}

\subsection{The structure of the errors}

Figure~\ref{fig:confusion} shows the row-normalised confusion matrices. The
errors are not spread arbitrarily over the class set: they concentrate in the
band adjacent to the diagonal, so a misassigned pixel usually lands on a
neighbouring interval of the scale axis rather than on an unrelated class. That
is the behaviour of a model estimating a scale and quantising it, and it is the
most direct support this study offers for the formulation of
Section~\ref{sec:intro}.

The two ends of the axis are effectively solved, open water and vast floe being
recovered at 98--99\%, while cake, small and medium floe exchange pixels
heavily. In the baseline, medium floe is recalled at 40.4\% with 38.1\% of its
pixels assigned to small floe: very nearly as many medium-floe pixels are read
as small floe as are read correctly. For big floe the baseline recovers 65.8\%,
sending 16.7\% to medium floe, 15.5\% to vast floe and 1.5\% to small floe. This
sharpens the description given in Section~\ref{sec:intro}. The loss on big floe
is not fragmentation into small floe, which accounts for very little of it; it
is a roughly even division between reading a large floe as a finer class and
absorbing it into the class above. Both are failures to localise a scale, and
both are consistent with scale branches that decide without reference to one
another, but only the first is fragmentation.

\begin{figure}[t]
\centering
\includegraphics[width=\linewidth]{fig5_confusion.pdf}
\caption{Row-normalised confusion matrices on fold~0 (percentage of each true class's pixels).
Confusion is concentrated in the band adjacent to the diagonal: errors move a class to its
neighbour on the scale axis rather than to an arbitrary class. Open water and vast floe are
essentially solved (98--99\%), while cake, small and medium floe exchange pixels heavily.}
\label{fig:confusion}
\end{figure}

\subsection{What does not improve}

Medium floe stands at 0.229 in the baseline and 0.229 in the full model. No
configuration on the main chain moves it, and the gate configuration that
improves the hard classes most, (b), moves it by $+0.002$. The bottleneck
identified as motivation~B in Section~\ref{sec:intro} --- medium floe
generalising poorly across scenes --- is therefore diagnosed by this work and
not resolved by it. We state this plainly rather than absorbing medium floe into
the hard-ice average, where a $+0.018$ gain on small floe would otherwise mask
it.

Cake ice remains very low throughout, rising only from 0.081 to 0.088. As set
out in Section~\ref{sec:data}, cake ice occupies few pixels and is heavily
confused with small floe; it is reported for completeness and is not a target of
this work.

\subsection{Additional probes}

We also examined two further modules, and found neither worth retaining.

We examined adding an adversarial scene-alignment objective on top of (b),
configuration (d), on the reasoning that if the failure is one of cross-scene
generalization then aligning scene representations should help. It does not:
macro mIoU falls to 0.4815, nearly 0.02 below (b), with lower IoU in five of the
six classes, including open water (0.951) and big floe (0.587). The single
exception is medium floe, which reaches 0.251 --- the highest value anywhere in
the table. We record that observation without building on it. It comes from one
fold, in a configuration that is worse almost everywhere else, and we have no
established account of the mechanism that would produce it.
[TODO: readouts for (d), if they bear on the interpretation.]

We also examined a polarimetric bias attention module on top of (c),
configuration (e). Its macro mIoU is 0.5103 against 0.5109 for (c): no net
change. What it does instead is move error between classes. Big floe rises to
0.700, the highest in the table, while medium floe falls to 0.206 and the
hard-ice diagnostic to 0.2193, both the lowest in the table. A module that adds
parameters and redistributes error without a net gain, in a study whose stated
concern is the hard classes, has nothing to recommend it, and it was removed.

Neither probe is part of PGSG. We report both because the design of
Section~\ref{sec:readouts} was intended to make negative outcomes decidable
rather than merely unpublished.

\subsection{Component-level readouts}

Figure~\ref{fig:readouts} reports the readouts of Section~\ref{sec:readouts}
over training. They answer whether a module is used, which is a different
question from whether it helps, and here the two answers do not agree.

The gate starts from $g \equiv 1$ on all four scales, as the zero-initialisation
requires, and ends with mean gate values of 1.14, 1.13, 1.02 and 0.71 from the
finest scale to the coarsest. Its selection entropy falls from 1.386 to 1.316
against a uniform upper bound of 1.386. The gate is therefore doing something
systematic, and what it does is a mild, consistent re-weighting away from the
coarsest scale rather than a per-pixel choice among scales; an entropy 0.07
below the uniform bound is a long way from selective allocation, and we quote
the bound alongside the value because the vertical axis of the panel is
magnified.

The cross-scale path is used more heavily. The residual added by the scale-axis
transformer grows from zero to 0.493 of the norm of the features it is added to,
which is consistent with that module producing the largest single-class change
in the study. Inside the layer, the evidence token draws 0.235 of the attention
mass against a uniform share of 0.200, so the evidence is attended to somewhat
more than a neutral token would be.

The remaining readout runs against this paper's account of its own method. The
share of the gate's input activation contributed by the evidence channels falls
from 0.394 early in training to 0.045 at the end. On that measure the trained
gate is driven almost entirely by the main features, and the evidence
conditioning from which the method takes its name has largely faded out of the
gate by the time training finishes --- even though the evidence remains in use
inside the scale-axis transformer. We report it because the readouts were
introduced precisely so that an outcome of this kind would be visible rather
than invisible, and we return to it in Section~\ref{sec:discussion}.

\begin{figure}[t]
\centering
\includegraphics[width=\linewidth]{fig6_readouts.pdf}
\caption{Training-time readouts of the added modules (logarithmic axis; readouts are logged at
geometrically spaced intervals). The gate starts from $g \equiv 1$, confirming the
zero-initialisation, and settles into a mild preference against the coarsest scale. The
cross-scale residual grows from zero to about half the norm of the scale features, so the
communication path is used. The share of the gate's input activation contributed by the
evidence decreases over training.}
\label{fig:readouts}
\end{figure}

\subsection{Qualitative comparison}

Figure~\ref{fig:qualitative} shows two fold-0 validation scenes. In the first, a
marginal-ice-zone scene, the reference has a band of medium floe near the coast
grading into small floe further out; the baseline reads most of the band as
small floe, and PGSG recovers a large part of it, with pixel agreement over the
displayed window rising from 0.700 to 0.808. The second is a failure case and is
shown as one. The reference there is uniformly big floe, and both models read it
as vast floe interspersed with patches of small floe; PGSG produces fewer such
patches, 0.170 agreement against 0.093, but both remain far from the reference.
The second scene is a useful corrective to the aggregate numbers: a class-mean
IoU of 0.676 on big floe is compatible with individual scenes on which the model
is wrong almost everywhere.

\begin{figure}[t]
\centering
\includegraphics[width=\linewidth]{fig7a_qualitative_hardclasses.pdf}\\[2pt]
\includegraphics[width=\linewidth]{fig7b_qualitative_bigfloe.pdf}
\caption{Qualitative comparison on two fold-0 validation scenes. Top: a marginal-ice-zone scene
in which the baseline reads most of the medium-floe band as small floe and PGSG recovers a
large part of it (pixel agreement over the displayed window 0.700 versus 0.808). Bottom: a
failure case in which the reference is uniformly big floe and both models read it as vast floe
interspersed with patches of small floe; PGSG produces fewer such patches, but both remain far
from the reference (0.093 versus 0.170).}
\label{fig:qualitative}
\end{figure}

\subsection{Scope of the evidence}

The differences discussed above are of order 0.01 or smaller in macro mIoU and
are measured on a single fold. They are consistent in sign and in their
class-wise structure with the mechanisms that produce them, which is the main
reason we regard them as more than noise; but a single fold cannot establish
that, and we do not claim it does. Table~[TODO] reports the paired five-fold
comparison between PGSG and the baseline, which is the measurement that settles
the question. Until then the reader should treat the fold-0 ordering as
indicative.


%% [TODO Table 4: five-fold paired comparison, PGSG vs baseline. Report the
%% per-fold difference, then the mean and standard error over the five paired
%% differences. Until this exists, every comparison in the paper stays inside
%% fold 0. Once it exists, fold-0 values and five-fold aggregate values still
%% must not be placed side by side as if comparable: they are different
%% evaluation sets.]

\section{Discussion}
\label{sec:discussion}

\subsection{Two failure modes, two mechanisms, and the tension between them}

The ablation separates cleanly along the lines the problem formulation
predicted. Evidence conditioning acts on the hard classes and costs accuracy on
the large ones; cross-scale communication acts on the large classes and returns
a little of the hard-class gain. Each mechanism moves the classes its motivation
named, which is the strongest support the fold-0 data can offer for the
formulation of Section~\ref{sec:intro}.

The interesting part is that the two are not simply additive. Configuration (b)
reaches a higher hard-ice value than the full model, and the full model reaches
a higher macro mIoU than either. One reading is that the gate and the scale-axis
layer compete for the same decision budget: once coarse-scale context is allowed
to reach the fine branch, pixels that the gate had assigned to a fine scale on
textural evidence are reassigned to a coarser one, which is exactly the
behaviour that repairs large-floe fragmentation and exactly the behaviour that
would cost a few hard-class pixels. We offer this as an interpretation, not as a
result. The readouts rule out the simplest alternative --- that one of the two
modules is inert --- since the gate departs measurably from uniform and the
cross-scale residual reaches roughly half the norm of the features it modifies.
What they do not show is that the reassignment described above is what the two
modules are trading, and we have not designed a probe that would. Should it hold, the practical
consequence is that the two mechanisms are better understood as a choice about
where to spend the scale budget than as independent improvements to be stacked.

\subsection{Why medium floe may lie outside the reach of the present evidence}

Medium floe does not move. Our design parameters suggest one candidate
explanation, which we set out because it is checkable rather than because the
present experiments establish it.

The variogram evidence is computed on stride-4 features, so its smallest lag is
320~m on the ground. The small--medium class boundary is at 100~m, and the whole
small-floe class spans 20--100~m. Every lag in the profile is therefore coarser
than the distinction the classifier most needs to draw. The profile can speak to
that distinction only indirectly, through how a field of such floes aggregates
at scales of 320~m and above. That the gate nonetheless improves small floe by
0.018 suggests the indirect route carries some information; that medium floe
does not move at all suggests it carries less where the two classes must be
separated from each other.

This is not a defect that a smaller lag would obviously repair, which is why we
do not present it as a fix. The lower bound on the lag was chosen for a reason
(Section~\ref{sec:variogram}): speckle is multiplicative noise concentrated at
the finest scales, so a lag near one pixel measures speckle rather than texture.
The design is caught between a noise floor from below and a class boundary from
above, and the interval between them is narrow at 80~m pixel spacing. Whether
finer-resolution imagery would widen it is a question this work cannot answer,
since it uses EW-mode data throughout.

Figure~\ref{fig:variogram} bears on this directly. The normalised class medians
do order monotonically with floe size, so the profile carries scale information.
But the interquartile ranges of the four classes overlap over most of the lag
range, and the big-floe and vast-floe medians very nearly coincide. Two of the
six classes are thus close to indistinguishable in this evidence at the lags
available, and the separation that does exist between small and medium floe is a
separation of medians inside heavily overlapping distributions. The evidence is
genuine and its invariance is exact, but on the present grid it is not sharp.

Two further possibilities deserve listing without being ranked. Supervision is
polygon-constant (Section~\ref{sec:data}), so a medium-floe polygon labels a
region whose interior may contain a range of floe sizes; the ceiling on medium
floe may partly be a property of the labels rather than of the model. And the
archive is geographically concentrated, so the held-out scenes for medium floe
may differ from the training scenes in ways that no amount of radiometric
invariance addresses.

\subsection{What the negative probes suggest}

Adversarial scene alignment is the standard response to a diagnosis of
cross-scene generalization failure, and here it degrades almost every class. One
explanation consistent with the dataset is that scene identity is not a pure
nuisance variable in this archive: scenes differ by region and season, and those
differences carry genuine information about which floe size classes are present
at all. Removing what distinguishes scenes therefore removes signal along with
noise. If that is right, the useful move is the one PGSG makes --- neutralising
a specific, named nuisance factor, namely absolute radiometric level, by
constructing evidence that is invariant to it --- rather than aligning scene
representations wholesale. We advance this as an account of the observation, not
as something the experiment demonstrates; distinguishing the two would require
probes that separate radiometric from compositional differences between scenes.

The polarimetric bias attention probe produced no net change while moving error
between classes. A plausible reading is that polarimetric information already
enters the network twice --- as an input channel (Section~\ref{sec:channels})
and as evidence (Section~\ref{sec:polar}) --- so a third pathway adds capacity
without adding information. Again, the readouts are what would settle this.

\subsection{The evidence stops driving the gate}

The most uncomfortable result in this study is a readout rather than a metric.
The share of the gate's input activation contributed by the evidence decays from
0.394 to 0.045 over training, while inside the scale-axis transformer the
evidence token holds slightly above a uniform share of the attention mass. Taken
at face value, the gate --- the module through which evidence conditioning is
meant to act --- ends up driven almost entirely by the main features, which are
the very features the evidence was introduced to bypass.

Two readings are available and we cannot separate them with what we have. The
first is that the gate's main-feature input already carries the information the
evidence supplies, having learned it during training, in which case the evidence
is redundant at that site and the hard-class gain of configuration (b) has some
other origin than the one the method's name asserts. The second is that the
decay is an artefact of optimisation: 96 evidence channels concatenated with
1024 feature channels are a small share of the gate input to begin with, and
nothing in the objective rewards keeping them influential, so the pathway may
simply be abandoned by gradient descent regardless of its merit.

The two readings call for different remedies --- the first would mean the gate
is the wrong site for conditioning, the second that the conditioning must be
protected, for instance by giving the evidence its own normalised pathway into
the gate or by an auxiliary term that penalises its disuse. Distinguishing them
is the clearest single target we can name for further work. We raise it here
rather than leave it to a reader to notice, because the readouts that expose it
were part of the design.

\subsection{Beyond this task}

The principle that organises the method --- construct evidence that is invariant
to a named nuisance factor, and let it condition the decision mechanism rather
than join the representation --- is not specific to floe size. It applies
wherever a label set is defined by a physical quantity that the sensor measures
only indirectly, and where the nuisance factors are known and nameable. What the
present study offers on that point is modest but concrete: an instantiation in
which the invariance is exact rather than learned, the conditioning path is
architecturally separated from the main path, and the zero-initialisation
discipline makes the contribution of each part separately decidable.

\section{Limitations}
\label{sec:limitations}

\paragraph{Evaluation.} The complete ablation exists on fold~0 only. The
differences reported are of order 0.01 or smaller in macro mIoU, which a single
fold cannot distinguish from run-to-run variation with confidence; the paired
five-fold comparison is the measurement that would, and it is pending
(Table~[TODO]). Seed variability has been estimated only from two runs of the
baseline, which differ by 0.005; that is enough to establish the scale of the
radiometric effect discussed in Section~\ref{sec:intro} but not to serve as a
variance estimate for the ablation. And because
no public benchmark exists for per-pixel floe size segmentation at EW
resolution, there is no external baseline to place these numbers against; we
have likewise not run a backbone comparison under the present cross-validation
protocol, so the choice of encoder is not empirically defended here.

\paragraph{Data.} The archive is geographically concentrated: 398 of 417 scenes
come from East Greenland and 19 from the Canadian Arctic. Grouped
cross-validation prevents a scene and its neighbouring frames from straddling
the train/validation boundary, but it cannot remove a correlation between class
and region that is built into the sampling itself.
%% [TODO] Once the per-region class pixel breakdown is available, quantify that
%% correlation here. Do not restate the earlier working assumption that big and
%% vast floe come mainly from the Canadian Arctic; it is unsupported.
Seasonal coverage is uneven, with autumn represented by 12 scenes, so nothing
in this work should be read as characterising autumn conditions. The dataset was
assembled for this study and is not a public release, which limits direct
reproduction.

\paragraph{Labels.} Supervision is polygon-constant: each ice-chart polygon
contributes a single floe size label to every pixel it covers. This is regional
weak supervision, it contains no individual floe boundaries, and its accuracy
depends on the dominant-form thresholds of Section~\ref{sec:data}, which we
deliberately loosened to retain marginal-ice-zone recall. Cake ice in particular
occupies very few pixels and is heavily confused with small floe; we report it
but do not treat it as an optimisation target, and its IoU should not be read as
a meaningful measurement.

\paragraph{Method.} The polarimetric prior attenuates but does not remove
incidence-angle dependence, since the angular slope of backscatter is itself
polarisation- and ice-type-dependent \citep{aldenhoff2020incidence}; a residual
angular term therefore survives the HH$-$HV difference. The variogram evidence
is contrast-invariant by construction, but its invariance is to affine
transformations of the feature field, which is a model of radiometric variation
and not radiometric variation itself; multiplicative speckle, incidence-angle
gradients across a swath, and thermal noise residuals are not affine and are not
covered by the guarantee. The radiometric sensitivity probe of
Section~\ref{sec:intro} applies a synthetic uniform shift, which is likewise a
simplification of how scenes actually differ.

\paragraph{Scope of the claim.} The central bottleneck this work set out to
address --- cross-scene generalization on medium floe --- is diagnosed and not
solved. The overall gain is consistent across classes in the sense that no class
is materially degraded, but it is moderate, and it is concentrated on large
floes rather than on the classes that motivated the study. The mechanistic
account is incomplete as well: the readouts show the evidence contribution to
the gate decaying over training (Section~\ref{sec:discussion}), so attributing
the hard-class gain of configuration (b) to evidence conditioning is less secure
than the name of the method suggests.

\section{Conclusions}
\label{sec:conclusions}

%% [TODO]

\section*{Acknowledgements}

%% [TODO] Includes the acknowledgement wording required for the ice-chart
%% sources.

\appendix

\section{[TODO: appendix title]}

\bibliographystyle{elsarticle-harv}
\bibliography{references}

\end{document}